\chapter{Reproducibility and Evidence Record}
\label{app:reproducibility}
\section{Source and Software}
The numerical inputs are from repository data revision
\texttt{\EvidenceRevision}. Figure/data SHA-256 hashes are written to
\path{generated/evidence_manifest.json} when the document is built. The
archived numerical run retains its own source hashes, configuration, model
checksums and stage records. A later report revision does not change those
underlying predictions.

\begin{table}[!htbp]
\centering\small
\caption{Software versions recorded with the numerical evidence.}
\label{tab:software}
\begin{tabularx}{\textwidth}{@{}YY@{}}
\toprule Package & Recorded version\\\midrule
\TableRows{software}
\bottomrule
\end{tabularx}
\end{table}
The scientific code was exercised under Python 3.12. These are numerical-run
versions, not a claim that a PDF build retrained the models. The proposal
builder itself uses the Python standard library and a standard LaTeX
installation; PyTorch is not needed simply to compile this document.

\section{Building the Proposal}
From a checkout of the repository, run:
\begin{lstlisting}
python proposal/build_proposal.py --bundle
\end{lstlisting}
The builder regenerates tables and highlighted values, checks required
inputs, runs latexmk/pdfLaTeX to resolve references, and writes
\path{proposal/Particle_PINN_Proposal.pdf}. The optional portable source ZIP
includes the numerical inputs and all required figure PDFs. The original
user-supplied template is preserved under \path{proposal/template/}.
For a new numerical run, use:
\begin{lstlisting}
python proposal/build_proposal.py --results results/NEW_RUN
\end{lstlisting}
The schema and required benchmarks must match; missing data trigger an
explicit error rather than numerical substitutions. Generated tables follow
the selected data. The interpretive prose is a reviewed narrative and must
also be checked when the scientific case or findings change.

\section{Inspecting and Reproducing the Numerical Study}
The following commands inspect the existing tests and re-evaluate saved
weights using the repository workflow. Install the scientific dependencies
as described in the project README before running them.
\begin{lstlisting}
python -m unittest discover -s tests -v
python run_particle_pinn.py --preset standard --seeds 2 \
  --ns 60 --nn 12 --weights-from results/final_study \
  --out results/proposal_replay
python assess_results.py results/proposal_replay
\end{lstlisting}
The replay writes a fresh output directory. It evaluates existing weights;
it does not reproduce the original training time. Model provenance paths
inside the old metadata are historical records; the required weights are
present under \path{results/final_study/models/}. New training is performed
with the documented training options rather than the replay command.

For an interior fluid marker with an arbitrary admissible start,
\begin{lstlisting}
python trace_path.py --run results/final_study \
  --start 0.35 0.35 --tracer --out results/interior_tracer
\end{lstlisting}
The equivalent particle command omits \texttt{--tracer}; an initial velocity
can be supplied explicitly. A new geometry or Reynolds number requires
retraining under the present per-case design.

\section{Complete Numerical Screening Record}
\Cref{tab:quality} is generated from the stored criterion flags. These flags
retain distinctions between execution, field accuracy and transport accuracy.
They are numerical screens, not an institutional approval or experimental
validation statement.
\begin{longtable}{@{}p{.70\textwidth}p{.23\textwidth}@{}}
\caption{Recorded numerical screening outcomes.}\label{tab:quality}\\
\toprule Check & Outcome\\\midrule
\endfirsthead
\toprule Check & Outcome\\\midrule
\endhead
\TableRows{quality}
\bottomrule
\end{longtable}

\section{Availability and Review}
The repository contains code, weights, raw numerical outputs, vector figures,
PNG images, a GIF and result reports. The animation is available as
\href{https://github.com/mohrezsayeh90-ai/basic_particle_problem/blob/main/results/final_study/figures/17_particle_motion.gif}{the particle-motion GIF}.
A PDF figure cannot encode that animation, so the underlying paths are also
shown as static figures and retained in CSV form. Access follows the
repository's current GitHub permissions. The supplied thesis template is
retained as a user-provided formatting source; no broader reuse license or
institutional approval is inferred from it.
